% Generated draft: requires mathematical and source review.
\documentclass[11pt]{article}
\usepackage[margin=25mm]{geometry}
\usepackage{fontspec}
\setmainfont{Latin Modern Roman}
\usepackage{amsmath,amssymb,mathtools,unicode-math}
\setmathfont{Latin Modern Math}
\usepackage{enumitem,needspace,xurl,hyperref}
\setlength{\parindent}{0pt}
\setlength{\parskip}{5pt}
\setlength{\emergencystretch}{3em}
\newcommand{\N}{\mathbb{N}}
\newcommand{\Z}{\mathbb{Z}}
\newcommand{\Q}{\mathbb{Q}}
\newcommand{\R}{\mathbb{R}}
\newcommand{\C}{\mathbb{C}}
\newcommand{\F}{\mathbb{F}}
\newcommand{\A}{\mathbb{A}}
\newcommand{\PP}{\mathbb{P}}
\newcommand{\E}{\mathbb{E}}
\newcommand{\eps}{\varepsilon}
\newcommand{\dd}{\,\mathrm{d}}
\begin{document}
\section*{UnsolvedMath ID: 701449}
\textbf{Source problem code:} catalogue entry
\subsection{Definitions and mathematical statement}
Let \(G\) be a finite group. Under the standard interpretation of characters as ordinary complex characters, a character is the trace function of a finite-dimensional complex representation of \(G\). Write \(\operatorname{Irr}_{\C}(G)\) for the irreducible such characters. The degree of \(\chi\in\operatorname{Irr}_{\C}(G)\) is \(\chi(1_G)\).

A normal subset \(M\subseteq G\) means a conjugation-invariant subset: \(gMg^{-1}=M\) for every \(g\in G\). It need not be a subgroup. For the formulation below, two functions on the same domain are called proportional if one is a nonzero complex scalar multiple of the other.

Characters \(\varphi,\psi\) are \emph{semiproportional} if they are not proportional on \(G\), but there exist a normal subset \(M\subseteq G\) and scalars \(a,b\in\C\setminus\{0\}\) such that
\[
\varphi(g)=
\begin{cases}
a\psi(g),&g\in M,\\
b\psi(g),&g\in G\setminus M.
\end{cases}
\]
No assumption that \(M\) is a subgroup, or that either part contains the identity, is imposed.

\textbf{Conjecture (V. A. Belonogov, Kourovka Notebook Problem 18.20).} For every finite group \(G\) and every pair of semiproportional characters \(\varphi,\psi\in\operatorname{Irr}_{\C}(G)\),
\[
\varphi(1_G)=\psi(1_G).
\]

\textbf{Unresolved conventions.} The supplied record does not explicitly specify the character field or whether proportionality permits a zero scalar. The displayed formulation uses ordinary complex characters and nonzero scalar proportionality; these conventions require confirmation against the cited source. The registry labels the conjecture open, but that status has not been independently verified.
\subsection{Short English statement}
Belonogov's conjecture says that two irreducible characters of a finite group have equal degrees whenever they are not globally proportional but become proportional separately on a conjugation-invariant subset and its complement. The complex-character and proportionality conventions stated above remain subject to source confirmation.
\subsection{Sources}
\begin{itemize}
\item Local UnsolvedMath registry, permanent ID \(701449\), exact source problem code \texttt{catalogue entry}, title \emph{Kourovka Notebook Problem 18.20}. The supplied record attributes the problem to V. A. Belonogov, gives the proposed year as 2014, and labels its status \texttt{open}. Its supplied creation and update timestamps are both \texttt{2026-10-16T14:00:00Z}.
\item Supplied source attribution: \emph{The Kourovka Notebook (Unsolved Problems in Group Theory)}, 18th issue (2014), Problem 18.20; the record also cites the 21st issue edition (2026), \texttt{arXiv:1401.0300v48}. The supplied collection description names E. I. Khukhro and V. D. Mazurov as editors of Issues 1--21 (1965--2026). These references and attribution are supplied evidence, not independently checked sources.
\end{itemize}
\end{document}
